\Problem{Triple Dobles}


\Statement

A cada partit de bàsquet es publiquen les estadístiques de cada jugador. Aquestes estadístiques recullen 5 categories: punts, rebots, assistències, taps i recuperacions. Es diu que un jugador ha fet un \emph{triple-doble} quan té un valor mínim de 10 a almenys 3 de les 5 categories anteriors. Un tipus especial de triple-doble és el \emph{doble triple-doble}. Això es produeix quan té un valor mínim de 20 a almenys 3 de les 5 categories. Per exemple:

\begin{itemize}
\item 20 10 5 15 30 i 13 10 5 15 30 : són triple-dobles
\item 5 20 10 20 25 i 15 20 10 20 25: són doble triple-dobles
\end{itemize}

Donades les estadístiques de les jugadores de la LFB en els partits ja disputats, les volem llistar ordenades segons el número de doble triple-dobles i triple-dobles que han fet.

És obligatori implementar i utilitzar el següent:

\begin{verbatim}
@dataclass
class Player:
    name: str
    dtd: int      # número de doble triple-dobles
    td: int       # número de triple-dobles (que no són doble triple-doble)

def read_player(p: int) -> Player:
    '''
    Pre: p >= 0; a l'entrada hi ha un nom seguit de les seves estadístiques a p partits
    Post: retorna la informació d'una jugadora segons les dades que hi havia a l'entrada
    '''
    ...

\end{verbatim}

\Input

L'entrada comença amb un enter $p \ge 0$ que representa el número de partits jugats. A continuació, segueix un enter $n >= 0$ i una seqüència amb la informació de $n$ jugadores. Per cada jugadora, ens donen el seu nom i les seves estadístiques al llarg dels $p$ partits jugats. Per cada partit, apareixen els punts, rebots, assistències, taps i recuperacions.

\Output

Volem veure la informació de les jugadores (nom, número de dobes triples-dobles i número de triples-dobles) ordenats segons els següents criteris: primer les que tinguin més dobles triple-dobles, en cas d'empat, les que tinguin més triples-dobles (que no siguin dobles triple-dobles) i, si persisteix l'empat, per ordre lexicogràfic del seu nom.


\Sample




